\documentclass[10pt,a4paper]{amsart}
\usepackage[margin=19mm]{geometry}
\usepackage{amsmath,amssymb,mathtools}
\usepackage[scr=rsfs,cal=euler]{mathalfa}
\usepackage{xfrac}
\usepackage[unicode,hidelinks,bookmarks=false]{hyperref}
\newcommand{\ie}{i.e.\ }
\newcommand{\cf}{cf.\ }
\newcommand{\resp}{respectively\ }
\newcommand{\Subsection}{\subsection}
\providecommand{\nobreakdash}{\nobreak\hbox{-}}
\setlength{\emergencystretch}{2em}
\multlinegap0pt
\usepackage{mathtools}
\usepackage{amssymb}
\usepackage{xcolor}
\usepackage{tikz}
\usepackage[shortlabels]{enumitem}
\newtheorem{lemma}{Lemma}
\newtheorem{theorem}{Theorem}
\usepackage{cleveref}
\crefname{lemma}{Lemma}{Lemmas}
\crefname{theorem}{Theorem}{Theorems}
\def\E{\mathcal{E}}
\def\F{\mathcal{F}}
\def\T{\mathcal{T}}
\def\pw{\mathrm{pw}}
\title{A hybrid high-order method for the biharmonic problem}
\author{Yizhou Liang, Ngoc Tien Tran}
\date{Source: arXiv:2504.16608v3 (CC BY 4.0)}
\begin{document}
\maketitle

\noindent\emph{Source and licence.} A hybrid high-order method for the biharmonic problem. Version arXiv:2504.16608v3 (CC BY 4.0), distributed by arXiv under the Creative Commons Attribution 4.0 International licence, http://creativecommons.org/licenses/by/4.0/, as recorded in the arXiv licence metadata for this exact version and shown on the abstract page https://arxiv.org/abs/2504.16608v3. Full licence notice: \texttt{licenses/arxiv-2504.16608-CC-BY-4.0.txt}.\par\smallskip

\noindent\textit{Editorial scope.} Counted unit: the article's theorem thm:L2 (source0.tex lines 573--578). The article's complete original proof of it is delivered with it (source0.tex lines 580--625, one \texttt{proof} environment of 46 lines). No mathematics is written, repaired or re-derived: the statement and the proof are the article's own lines, in the article's own notation, and no retained line is substituted or re-flowed.  Retained uncounted as the source-local context the proof needs: lines 189--191; lines 236--241; lines 255--260; lines 300--302; lines 359--361; lines 496--500; lines 504--507; lines 513--518; lines 555--557; lines 569--571.  Nothing was omitted from any retained span, so the package carries no omission record.  No retained line contains a citation, so the wrapper carries no bibliography and no bibliography entry.  No retained line contains a figure environment, an image-inclusion command or drawing code, so no image is delivered and the package declares no asset dependency.  Editorial formatting only, in addition: an AMS \texttt{amsart} preamble that restates the article's own declarations and macros used by the retained lines, namely the environments \texttt{lemma}, \texttt{theorem} on separate counters, together with the standard packages those lines need.  The retained environments print in this document as Lemma 1, Theorem 1 in order of appearance, whereas the article prints the same items with the numbering of its own full document.  The complete native source of the article as distributed in the arXiv e-print for this exact version is bundled unchanged as \texttt{sources/176966/source0.tex}.
\begin{align}\label{def:source-problem}
    (\nabla^2 u, \nabla^2 v)_{\Omega} = (f,v)_{\Omega} \quad\text{for any } v \in H^2_0(\Omega).
\end{align}

\begin{lemma}[commuting]\label{lem:commuting}
    Any $v \in H^2(T)$ satisfies the $L^2$ orthogonality $\nabla^2 (v - R_T I_T v) \perp \nabla^2 P_{k+2}(T)$. In particular,
    \begin{align}\label{ineq:RI-best-appr}
        \|\nabla^2 (v - R_T I_T v)\|_T = \min_{p \in P_{k+2}(T)} \|\nabla^2 (v - p)\|_T.
    \end{align}
\end{lemma}

\begin{lemma}[optimality of $s_T$]\label{lem:best-appr-stab}
    Let a simplex $T$ be given. Any $v \in H^2(T)$ satisfies
    \begin{align*}
        \sqrt{s_T(I_T v, I_T v)} \lesssim \min_{p \in P_{k+2}(\T)} \|\nabla^2(v - p)\|_T.
    \end{align*}
\end{lemma}

\begin{align}\label{def:discrete_problem}
    a_h(u_h,v_h) = (f,v_\T)_\Omega \quad\text{for any } v_h = (v_\T, v_\F, \beta_\F, v_\E) \in V_h
\end{align}

\begin{align}\label{eq:commuting-G}
	\nabla_\pw^2(v - G_h v) \perp \nabla_\pw^2 P_{k+2}(\T)
\end{align}

	\begin{equation}\label{ineq:quasi-best-J-R}
    \begin{aligned}
        \|h_\T^{-2}(J_h - R_h) I_h v\|_\Omega &+ \|h_\T^{-1}\nabla_\pw(J_h - R_h) I_h v\|_\Omega + \|\nabla_\pw^2(J_h - R_h) I_h v\|_\Omega \\ & \lesssim \|\nabla_\pw^2(v - G_h v)\|_\Omega = \min_{p \in P_{k+2}(\T)} \|\nabla_\pw^2(v - p)\|_\Omega.
    \end{aligned}
	\end{equation}

	\begin{align}\label{ineq:quasi-interpolation}
		\|h_\T^{-2}(v - J_h I_h v)\|_\Omega + \|h_\T^{-1}\nabla (v - J_h I_h v)\|_\Omega + \|\nabla^2(v - J_h I_h v)\|_\Omega&\nonumber\\
		\lesssim \|\nabla_\pw^2(v - G_h v)\|_\Omega \leq \|\nabla^2 v\|_\Omega&
	\end{align}

\begin{theorem}[a~priori]\label{thm:a-priori}
    The discrete solution $u_h \in V_h$ to \eqref{def:discrete_problem} satisfies
    \begin{equation}\label{ineq:a-priori-thm}
        \|I_h u - u_h\|_{h} + |u_h|_{s_h} \lesssim \min_{p \in P_{k+2}(\mathcal{T})}\|\nabla^2_\pw(u - p)\|_{\Omega} + \mathrm{osc}(f,\mathcal{T})
    \end{equation}
\end{theorem}

    \begin{align}\label{eq:orth-R-J}
    	\nabla^2_\mathrm{pw} (R_h e_h - J_h e_h) \perp \nabla^2_\mathrm{pw} P_{k+2}(\mathcal{T})
    \end{align}

\begin{align}\label{ineq:elliptic-reg}
	\|z\|_{H^{2+s}(\Omega)} \lesssim \|g\|_\Omega.
\end{align}

\begin{theorem}[$L^2$ error]\label{thm:L2}
	If $\ell \geq 2$, the discrete solution $u_h \in V_h$ to \eqref{def:discrete_problem} and $t \coloneqq \min\{s,\ell-1\}$ satisfy
	\begin{align*}
		\|\Pi_\T^\ell(u - u_\T)\|_\Omega \lesssim h^{t} (\min_{p \in P_{k+2}(\mathcal{T})}\|\nabla^2_\pw(u - p)\|_{\Omega} + \mathrm{osc}(f,\T)).
	\end{align*}
\end{theorem}

\begin{proof}
    We adopt the abbreviation $e_h = (e_\T, e_\F, \gamma_\F, e_\E) = I_h u - u_h \in V_h$ from above. Let $z\in V$ solve $\Delta^2 z = e_{\mathcal{T}}$ in $\Omega$, i.e.,
    \begin{equation*}
    	(\nabla^2 z,\nabla^2v)_{\Omega} = (e_\T, v)_{\Omega} \quad \text{for any } v \in V.
    \end{equation*}
    Since $\Pi_\T^\ell J_h e_h = e_\T$ from $I_h \circ J_h = \mathrm{Id}$ in $V_h$, this implies the split
    \begin{align}\label{eq:proof-L2-split}
    	\|e_\T\|_\Omega^2 &= (e_\T, J_h e_h)_{\Omega} = (\nabla^2 z,\nabla^2 J_h e_h)_{\Omega}\nonumber\\
    	& = (\nabla^2 z,\nabla^2_\pw R_h e_h)_{\Omega} + (\nabla^2 z, \nabla^2_\pw (J_h e_h - R_h e_h))_{\Omega}.
    \end{align}
    The $L^2$ orthogonality \eqref{eq:orth-R-J}, a Cauchy inequality, and the continuity of $J_h$ prove
    \begin{align}\label{ineq:proof-L2-1}
    	(\nabla^2 z, \nabla^2_\pw (J_h e_h - R_h e_h))_{\Omega} & = (\nabla^2_\pw (z- G_h z),\nabla^2_\pw J_h e_h)_{\Omega}\nonumber\\
    	&\lesssim \|\nabla^2_\pw (z - G_h z)\|_\Omega \|e_h\|_h.
    \end{align}
    Linear algebra shows the split
    \begin{align*}
    	(\nabla^2 z,\nabla^2_\pw &R_h e_h)_{\Omega} = (\nabla^2 z, \nabla^2_\pw G_h u)_\Omega - (\nabla^2_\pw G_h z, \nabla^2_\pw R_h u_h)_\Omega\\
    	&= (\nabla^2 z, \nabla^2 u)_\Omega - (\nabla^2_\pw G_h z, \nabla^2_\pw R_h u_h)_\Omega - (\nabla^2 z, \nabla^2_\pw(u - G_h u))_\Omega.
    \end{align*}
    The first term on the right-hand side is equal to
    \begin{align*}
    	 (\nabla^2_\pw(z - G_h z), \nabla^2 u)_\Omega + (\nabla_\pw^2(G_h z - J_h I_h z), \nabla^2 u)_\Omega + (\nabla^2 J_h I_h z, \nabla^2 u)_\Omega.
    \end{align*}
    This, the previously displayed formula, the variational formulations \eqref{def:source-problem} and \eqref{def:discrete_problem}, and $(\nabla^2 z, \nabla^2_\pw(u - G_h u))_\Omega = (\nabla^2_\pw(z - G_h z), \nabla^2 u)_\Omega$ from \eqref{eq:commuting-G} provide
    \begin{align}\label{eq:proof-L2-identity}
    	(\nabla^2 z,\nabla^2_\pw R_h e_h)_{\Omega} &= (f, J_h I_h z - \Pi_\T^\ell z)_\Omega\nonumber\\
    	&\qquad + s(I_h z, u_h) + (\nabla_\pw^2(G_h z - J_h I_h z), \nabla^2 u)_\Omega.
    \end{align}
    The $L^2$ orthogonality $J_h I_h z - \Pi_\T^\ell z \perp P_{\ell}(\T)$ from $I_h \circ J_h = \mathrm{Id}$ in $V_h$ and \eqref{ineq:quasi-interpolation} imply
    \begin{align*}
    	(f, J_h I_h z - \Pi_\T^\ell z)_\Omega &= ((1 - \Pi_\T^\ell)f, J_h I_h z-z)_\Omega + ((1 - \Pi_\T^\ell)f, (1 - \Pi_\T^\ell) z)_\Omega\\
    	&\lesssim \mathrm{osc}(f,\T)(\|\nabla_\pw^2(z - G_h z)\|_\Omega + \|h_\T^{-2}(1-\Pi_\T^\ell)z\|_\Omega).
    \end{align*}
    Hence, \eqref{eq:proof-L2-identity}, the $L^2$ orthogonality $\nabla_\pw^2(G_h z - J_h I_h z) \perp \nabla_\pw^2 P_{k+2}(\T)$ from \Cref{lem:commuting}, the Cauchy inequality, \Cref{lem:best-appr-stab}, and \eqref{ineq:quasi-best-J-R} show
    \begin{align*}
    	(\nabla^2 z,\nabla^2_\pw R_h e_h)_{\Omega} \lesssim (\mathrm{osc}(f,\T) + |u_h|_{s_h} + \|\nabla_\pw^2(u - G_h u)\|_\Omega)&\\
        \times (\|\nabla_\pw^2(z - G_h z)\|_\Omega + \|h_\T^{-2}(1-\Pi_\T^\ell)z\|_\Omega)&.
    \end{align*}
    The combination of this with \eqref{eq:proof-L2-split}--\eqref{ineq:proof-L2-1} results in
    \begin{align*}
    	\|e_\T\|^2_\Omega \lesssim (\mathrm{osc}(f,\T) + \|e_h\|_h + |u_h|_{s_h} + \|\nabla_\pw^2(u - G_h u)\|_\Omega)&\\
        \times (\|\nabla_\pw^2(z - G_h z)\|_\Omega + \|h_\T^{-2}(1-\Pi_\T^\ell)z\|_\Omega)&.
    \end{align*}
    The assertion then follows from \Cref{thm:a-priori} and the elliptic regularity \eqref{ineq:elliptic-reg}.
\end{proof}

\end{document}
